\documentclass[11pt]{article}
\usepackage[margin=1in]{geometry}
\usepackage{amsmath,amssymb,amsthm,mathtools,microtype}
\usepackage[hidelinks]{hyperref}
\hypersetup{pdftitle={The outer focal-pedal area product in odd elliptic billiards},pdfauthor={Alec Kriebel}}
\newtheorem{theorem}{Theorem}
\newtheorem{lemma}[theorem]{Lemma}
\newtheorem{corollary}[theorem]{Corollary}
\newcommand{\sn}{\operatorname{sn}}
\newcommand{\cn}{\operatorname{cn}}
\newcommand{\dn}{\operatorname{dn}}
\newcommand{\ii}{\mathrm{i}}
\newcommand{\Area}{\mathcal A}
\title{The outer focal-pedal area product\ in odd elliptic billiards}
\author{Alec Kriebel\\Independent researcher\\
\small\href{https://orcid.org/0009-0001-9320-500X}{ORCID: 0009-0001-9320-500X}}
\date{15 October 2026}
\begin{document}
\maketitle
\begin{abstract}
For a fixed odd-periodic billiard family between strictly nested confocal
ellipses, project each original focus onto the ordered side lines of the outer
tangent polygon. We prove that the product of the two signed pedal areas is
constant, for every coprime winding, including stars. This proves the
experimental assertion labelled $k_{605,a}$ in the 2020 arXiv invariant list
and $k_{606}$ in its published companion. The focal-foot map has paired
collinear residues; its cyclic area consequently has only two possible simple
poles on a reduced torus. Reflection and an imaginary antiperiod provide the
complementary zeros. The argument uses established elliptic-function methods.
We explain connections to generalized Landen identities and recent focal
inversion results; historical priority of this particular application remains
undetermined. Exact controls and separately labelled numerical diagnostics
accompany this unrefereed, extensively AI-assisted note.
\end{abstract}

\section{The assertion and its precise formulation}
For an ordered closed polygon $W=(W_0,\ldots,W_{N-1})$, set
\begin{equation}\label{eq:area}
 \Area(W)=\frac12\sum_{j=0}^{N-1}\det(W_j,W_{j+1}),\qquad W_N=W_0.
\end{equation}
This is signed shoelace area in trajectory order, also for a star.
Let $E:x^2/a^2+y^2/b^2=1$, $a>b>0$, have foci
$F_\pm=(\pm c,0)$, where $c^2=a^2-b^2$. Given a billiard polygon
$P=(P_j)$, let $\ell_j$ be the tangent to $E$ at $P_j$, and let
$V_j=\ell_j\cap\ell_{j+1}$. The outer polygon has side lines
$V_{j-1}V_j=\ell_j$. Write $Q_{\pm,j}$ for the orthogonal foot of
$F_\pm$ on this line, and $B_\pm=\Area((Q_{\pm,j})_j)$.

\begin{theorem}\label{thm:main}
Fix a strictly nested, nondegenerate confocal elliptical caustic inside $E$
supporting a billiard family of odd least period $N\geq3$. Then $B_+B_-$
is independent of the initial vertex. Every admissible coprime star winding
is included, with the order and signed-area convention above.
\end{theorem}

Reznik, Garcia and Koiller state this primed-area product experimentally in
\cite[Table 7, p.~9, $k_{605,a}$]{RGK}, and in the published edition
\cite[Table 7, p.~349, $k_{606}$]{RGKfinal}. The published $k_{605}$ concerns
the different, original-chord pedal polygons. The outer construction is in
Figure~1 and Sections~3.4 and~3.7 of those sources; Section~2 specifies signed
area. The assertion already occurs as $k_{604}$ in Table~5, p.~6 of
arXiv:2004.12497v1, submitted 26~April~2020. The source tables do not
explicitly define least-period, star or repetition conventions; these are
made precise in the theorem and corollary here. Hyperbolic, collapsed and
grazing caustics are outside the stated domain.

\section{Canonical coordinates and the focal foot}
Let the caustic semiaxes be $\alpha>\beta>0$, so that
$\alpha^2-\beta^2=c^2$. Put $k=c/\alpha$, $k'=\beta/\alpha$,
$K=K(k)$ and $K'=K(k')$. In this section Jacobi functions use
\emph{modulus} $k$, not parameter $k^2$. Stachel's canonical
parametrization \cite[Theorem 4.3, equation (4.9), p.~1614]{Stachel} gives,
after choosing orientation,
\begin{gather}\label{eq:canonical}
 P(w)=(-a\sn w,b\cn w),\quad P_j(w)=P(w+j\delta),\qquad
 a=\frac{\alpha\dn v}{\cn v},\quad b=\frac{\beta}{\cn v},\\
 \delta=2v=\frac{4K\tau}{N},\qquad
 \gcd(\tau,N)=1,\quad 0<\tau<N/2,\quad 0<v<K.\notag
\end{gather}
Only windings admitted by the fixed confocal pair are asserted. This is
uniform canonical phase, rather than uniform eccentric anomaly.

The tangent at $P(w)$ is $n(w)\cdot X=1$, with
$n(w)=(-\sn w/a,\cn w/b)$. Perpendicular projection gives
\begin{equation}\label{eq:foot}
 Q(w)=Q_+(w)=
 \left(\frac{a(c-a\sn w)}{a-c\sn w},
       \frac{ab\cn w}{a-c\sn w}\right).
\end{equation}
Indeed $n\cdot n=(a-cs)(a+cs)/(a^2b^2)$ for $s=\sn w$, and the
factor $a+cs$ cancels in $F_++(1-n\cdot F_+)n/(n\cdot n)$.
The real denominator is at least $a-c>0$. Direct substitution gives
$Q(w)\cdot Q(w)=a^2$, the classical auxiliary-circle property
(also recorded as $k_{301}$ in \cite{RGK}). All complex dot products below
are bilinear, not Hermitian.

The real phase parametrization is injective modulo $4K$. Consecutive
vertices are distinct and nonantipodal since $0<\delta<2K$, so their
tangents meet finitely. Three distinct tangents to a nondegenerate conic
cannot be concurrent; thus no outer side collapses. The ordered feet on
the outer sides are precisely \eqref{eq:foot} at $w+j\delta$, up to a
cyclic shift. Define
\begin{equation}\label{eq:trace}
 T(w)=\frac12\sum_{j=0}^{N-1}
 \det\bigl(Q(w+j\delta),Q(w+(j+1)\delta)\bigr).
\end{equation}
The real half-period identities give
\begin{equation}\label{eq:focus}
 Q_-(w)=-Q(w+2K),\qquad B_+(w)=T(w),\quad B_-(w)=T(w+2K).
\end{equation}
Central negation leaves signed area unchanged.

\section{The two possible simple poles of the area}
We use standard periods, addition and quarter-period identities
\cite[Sections 22.4 and 22.8]{DLMF}. Set $p=\ii K'$ and $r=K+p$.
At a common pole of $\sn$ and $\cn$, the numerator and denominator in
\eqref{eq:foot} have at most simple poles with nonzero leading denominator,
so the singularity is removable. All other possible poles solve
\[
 \sn w=\frac ac=\frac{\dn v}{k\cn v}.
\]
The quarter-period identities are
\begin{equation}\label{eq:quarter}
 \sn(r+z)=\frac{\dn z}{k\cn z},\quad
 \cn(r+z)=\frac{-\ii k'}{k\cn z},\quad
 \dn(r+z)=\frac{\ii k'\sn z}{\cn z}.
\end{equation}
On the $\sn$ torus with periods $4K,2p$, the two roots are
$r-v,r+v$. They are distinct, and the derivative
$-c\cn w\dn w$ of the denominator is nonzero there.
Since $\sn$ has degree two on this torus, the list is exhaustive.
The common torus for the vector $Q$ has periods $4K,4p$;
it also includes the two roots translated by $2p$.
In particular, $2p$ must not be treated as a period of the full vector.

At $r-v$ and $r+v$, the values of $\sn$ and $\cn$ are respectively
$a/c$ and $-\ii b/c$. The two numerator vectors in \eqref{eq:foot}
are therefore both $(-ab^2/c,-\ii ab^2/c)$. The denominator derivatives
have opposite signs by \eqref{eq:quarter}. The residue vectors of $Q$
are consequently opposite and collinear.

\begin{lemma}\label{lem:poles}
On $X=\mathbb C/(L\mathbb Z+4p\mathbb Z)$, $L=4K/N$, the function
$T$ has at most simple poles, only at $b_0+p,b_0+3p$, where $b_0=K-v$.
\end{lemma}
\begin{proof}
The poles of $Q$ at $r-v$ and $r+v$ differ by $\delta=2v$.
At a singular phase of the cyclic sum, exactly two vertices are singular,
at adjacent indices: coprimality makes all $N$ phases distinct modulo
$4K$. There are no additional copies from the other imaginary row.
Write the two adjacent expansions as
\[
 Q(r-v+\varepsilon)=R/\varepsilon+A+O(\varepsilon),\qquad
 Q(r+v+\varepsilon)=-R/\varepsilon+B+O(\varepsilon).
\]
The possible double coefficient of their determinant is
$\det(R,-R)=0$. Every other edge determinant has at most one singular
endpoint, hence has at most a simple pole. This also handles $N=3$,
where the other two incident edges share their regular vertex, and the
edge crossing the cyclic index cut: $N\delta=4K\tau$ is a true period.
The same argument applies in the translated row.

Cyclic relabelling gives period $\delta$ for $T$; it also has period
$4K$. Bezout's identity with $\gcd(\tau,N)=1$ gives period $L$.
The possible pole phases $r-v-j\delta$ and their $2p$ translates
then collapse to the two displayed points on $X$.
\end{proof}

\section{Reflection, complementary zeros and the product}
Let $J=\operatorname{diag}(1,-1)$. The Jacobi identities give
\[
 Q(w+2p)=JQ(w),\qquad Q(2K-w)=JQ(w).
\]
The first changes each determinant sign. In the second, reflection and
cyclic order reversal each contribute a minus sign, which cancel. Thus
\begin{equation}\label{eq:characters}
 T(w+2p)=-T(w),\qquad T(2K-w)=T(w).
\end{equation}
Since $2v=\delta$ is a period, $F(z)=T(b_0+z)$ is even.
It has periods $L,4p$, antiperiod $2p$, and by Lemma~\ref{lem:poles}
at most simple poles at $p,3p$.

At the regular point $z_0=p+L/2$, evenness and real periodicity give
$F(-p-L/2)=F(-p+L/2)=F(p+L/2)$, whereas the antiperiod gives
$F(-p+L/2)=-F(p+L/2)$. Hence $F(z_0)=0$.
The same holds at $3p+L/2$ by translation. These points are distinct
from $p,3p$ on $X$. Every possible pole of
\begin{equation}\label{eq:product}
 F(z)F(z+L/2)
\end{equation}
is cancelled by one of those zeros in the other factor, because all
possible poles are simple. The product is holomorphic on a compact
torus and therefore constant. This also covers the pole-free case
(then the antiperiod forces $F=0$). No division by a possibly zero
area is involved. This elementary complementary-divisor argument is a
classical elliptic-function mechanism, not a new general method.

For odd $N$, $2K=(N/2)L\equiv L/2\pmod{L\mathbb Z}$.
Equations~\eqref{eq:focus} and~\eqref{eq:product} prove
$B_+(w)B_-(w)=\mathrm{constant}$, establishing Theorem~\ref{thm:main}.
The argument used no restriction to winding one, so includes all
admitted primitive stars.

\begin{corollary}
An $s$-fold traversal of a fixed odd primitive family has areas $sB_\pm$
and product $s^2B_+B_-$. Orientation reversal changes both area signs
and preserves the product. For a circular boundary, the coincident-focus
pedal polygons are regular stars and their product is constant by rotation.
\end{corollary}
\begin{proof}
Each ordered edge determinant occurs exactly $s$ times in a repeated
list; reversal negates the sum. On a circle the feet from its center to
its tangents are the original regular-star vertices.
\end{proof}
For even \emph{primitive} periods the focus shift is a full reduced real
period, yielding equal focal areas rather than constancy of their square.
The even countercontrols below show that the odd hypothesis is essential.

\section{Related identities, current work and priority}
The auxiliary circle, canonical coordinates, Jacobi identities and compact
torus reasoning are established tools. Complex continuation and pole
removal occur in earlier Poncelet-area work
\cite{AST,Chavez} and bicentric-invariant work \cite{RGR}.
General cyclic elliptic-function identities and principal-part
reconstruction also predate this application \cite{KLS}.

There is a second exact proof route using the old arbitrary-order Landen
identity of Khare and Sukhatme \cite[equation (7), p.~923; Section 2.1]{KS}.
For clarity, the full geometric reduction, sign analysis for stars and
an independent Fourier-series derivation are in the accompanying source
archive, under \path{verification/audits/priority_classical_general_01/}
in \path{EXACT_BICENTRIC_DN_BRIDGE.md}. Briefly, the focal foot is the
restriction to $E$ of the projectivity
\[
 (x,y)\longmapsto
 \left(\frac{a^2(x+c)}{a^2+cx},\frac{a^2y}{a^2+cx}\right).
\]
It maps the caustic to an ellipse with a focus at the auxiliary-circle
center. Polarity $q\cdot R=a^2$ then gives a genuine bicentric polygon
whose contact polygon is $Q$. With $\lambda=a^2-\alpha^2=b^2-\beta^2$,
its incircle radius, circumcenter and circumradius are
\[
 r_b=a,\qquad d_b=-c\lambda/\beta^2,\qquad R_b=\alpha b^2/\beta^2.
\]
A local addition identity and cyclic reindexing give
$B_+/\Area(R)=r_b^2/(R_b^2-d_b^2)$ and
\[
 B_+(w)=\frac{r_b^3t_M}{R_b^2-d_b^2}
 \sum_{j=0}^{N-1}\dn(u+2jK_b\tau/N,m_b),
\]
where this display uses \emph{parameter}
$m_b=4k/(1+k)^2$, $K_b=K(\sqrt{m_b})=(1+k)K$,
$u=(1+k)(w+v-K)/2$, and $t_M^2=(R_b+|d_b|)^2-r_b^2$.
The focus switch is exactly $u\mapsto u+K_b$.
The published arbitrary-order identity reduces this sum to a scaled
$\dn$ with nome $q_b^N$; for odd $N$ the focus switch becomes its
real half-period. The classical half-period product proves the same
constant. The analytic identity is old; the complete geometric
application above was derived during this audit. Its derivability does
not establish its earlier publication or its historical novelty.

Recent work of Ferudun proves the odd product for unit inversions of the
\emph{original vertices} about the two foci \cite{F62}, and a related
outer-pedal area product for other parity/reference-point hypotheses
\cite{F14}. The latter includes parallel-residue cancellation, a close
methodological antecedent. These results warrant substantive credit.
An explicit comparison in the archive,
under \path{verification/audits/priority_author_updates_01/}
in \path{CHECKABLE_TRANSFER.md}, shows that $B_e=\gamma D_{-e}$ for a
phase-independent scalar, where $D_e$ is the original-vertex unit-inversion
area. The additional outer-foot pole bound of Lemma~\ref{lem:poles},
together with the inverse area's proved pole and antiperiod facts,
puts both traces in the same at-most-one-dimensional function space.
Matching residues proves proportionality without dividing through zeros.
That geometric application is not printed in \cite{F62}. Its inverse
theorem with this additional lemma also proves the present product.
Neither an inverse/pedal label mismatch nor a generic projective
correspondence alone settles priority or transports Euclidean area.

Two independent, bounded priority investigations, integrated with primary
source checks, verified no earlier printed general proof of this exact
outer focal-pedal assertion. This is not an exhaustive literature result.
Historical novelty and public proof priority remain undetermined; no
first-proof, new general elliptic-method or current worldwide-open claim
is made. The source archive records actual editions, passages, hashes,
coverage and access gaps, including the inaccessible full text of
Stachel's \emph{A Triple of Projective Billiards} and the incompletely
compared final journal edition of \emph{The Talented Mr. Inversive Triangle}.
Open author versions and selected related primary proofs were inspected.

\section{Checkable evidence and research status}
For $a^2=21$, $b^2=16$, $\lambda=336/25$, two genuine triangular phases are
\[
 H=((\sqrt{21},0),(-3\sqrt{21}/5,16/5),(-3\sqrt{21}/5,-16/5)),
\]
\[
 V=((0,4),(-21/5,-8/5),(21/5,-8/5)).
\]
These example inputs also occur in \cite[Section 6]{F62}; their reuse is
not a new parameter discovery. Direct exact reflection, common-caustic
tangency, actual tangent intersections and independent perpendicular feet
give
\[
 B_\pm(H)=\frac{49\sqrt{21}\pm7\sqrt5}{10},\qquad
 B_\pm(V)=\frac{112}{5},\qquad B_+B_-=\frac{12544}{25}.
\]
The original-chord pedal product is instead $7225344/390625$ in these
phases. The original-vertex unit-inverse product is $1/324$, and
$B_e=(2016/5)D_{-e}$ in both phases. These finite exact checks distinguish
the objects; the universal claims are proved above and in the cited
comparison derivations.

The portable archive preserves the original submission, independent
audit reports, exact symbolic and finite lattice controls and separately
labelled high-precision diagnostics. New checks include physical
reflection stepping without Jacobi coordinates, complex pole behavior,
near-circular and thin-caustic families, stars and even negative controls.
Numerical agreement has no interval certification and is not a substitute
for the all-period proof. The package gives pinned dependencies, hashes
and an isolated replay harness, with source-access and reproducibility
limits stated explicitly.

AI tools were used extensively in solving, source investigation, algebra,
computation, writing and independent internal adversarial checking.
The author is responsible for the claims. This is an unrefereed preprint;
no independent human peer review or formal proof-assistant verification
is represented. The proof, original submission and verification materials
are provided so that independent readers can check the result.

\begin{thebibliography}{99}\small
\bibitem{RGK} D. Reznik, R. Garcia and J. Koiller,
\emph{Eighty New Invariants of N-Periodics in the Elliptic Billiard},
arXiv:2004.12497v11 (2020). Earlier v1 (26 April 2020) is titled
\emph{Forty New Invariants of N-Periodics in the Elliptic Billiard}.
\url{https://arxiv.org/abs/2004.12497v11}.
\bibitem{RGKfinal} D. Reznik, R. Garcia and J. Koiller,
\emph{Fifty New Invariants of N-Periodics in the Elliptic Billiard},
Arnold Mathematical Journal \textbf{7} (2021), 341--355.
\href{https://doi.org/10.1007/s40598-021-00174-y}{doi:10.1007/s40598-021-00174-y}.
\bibitem{Stachel} H. Stachel, \emph{On the motion of billiards in ellipses},
European Journal of Mathematics \textbf{8} (2022), 1602--1622.
\href{https://doi.org/10.1007/s40879-021-00524-2}{doi:10.1007/s40879-021-00524-2}.
\bibitem{DLMF} F. W. J. Olver et al., eds., \emph{NIST Digital Library of
Mathematical Functions}, Sections 22.4, 22.7, 22.8 and 22.11.
\url{https://dlmf.nist.gov/22}.
\bibitem{AST} A. Akopyan, R. Schwartz and S. Tabachnikov,
\emph{Billiards in ellipses revisited}, European Journal of Mathematics
\textbf{8} (2022), 1313--1327. Author version:
\url{https://arxiv.org/abs/2001.02934v2}.
\bibitem{Chavez} A. Ch\'avez-Caliz, \emph{More About Areas and Centers of
Poncelet Polygons}, Arnold Mathematical Journal \textbf{7} (2021), 91--106.
\url{https://armj.math.stonybrook.edu/pdf-Springer-final/020-0154.pdf}.
\bibitem{RGR} P. Roitman, R. Garcia and D. Reznik,
\emph{New Invariants of Poncelet--Jacobi Bicentric Polygons},
Arnold Mathematical Journal \textbf{7} (2021), 619--637.
\href{https://doi.org/10.1007/s40598-021-00188-6}{doi:10.1007/s40598-021-00188-6}.
\bibitem{KLS} A. Khare, A. Lakshminarayan and U. Sukhatme,
\emph{Cyclic Identities Involving Jacobi Elliptic Functions. II},
Journal of Mathematical Physics \textbf{44} (2003), 1822--1841.
\url{https://arxiv.org/abs/math-ph/0207019v2}.
\bibitem{KS} A. Khare and U. Sukhatme,
\emph{Connecting Jacobi elliptic functions with different modulus parameters},
Pramana \textbf{63} (2004), 921--936.
\href{https://doi.org/10.1007/BF02704331}{doi:10.1007/BF02704331}.
\bibitem{F62} A. Ferudun, \emph{Products of the Two Focal Inversion Areas
in Elliptic Billiards}, unrefereed preprint, version 1.0, 2 October 2026.
\href{https://doi.org/10.5281/zenodo.23096994}{doi:10.5281/zenodo.23096994}.
\bibitem{F14} A. Ferudun, \emph{Outer-Pedal Area Products in Elliptic
Billiards}, unrefereed preprint, version 1.1, 1 October 2026.
\href{https://doi.org/10.5281/zenodo.23088516}{doi:10.5281/zenodo.23088516}.
\end{thebibliography}
\end{document}
